\section{Theoretical Background}
\label{sec:background}

This section provides the theoretical foundation for the proposed CAC-CycleGAN-WGP framework. We first review the fundamental principles of Generative Adversarial Networks (GANs), followed by the mechanism of Cycle-Consistent GAN (CycleGAN) for unpaired domain translation. Finally, we discuss the Wasserstein distance and Gradient Penalty (GP) used to ensure training stability and Lipschitz continuity.

\subsection{Generative Adversarial Networks (GAN)}
The GAN framework, introduced by Goodfellow et al. \cite{goodfellow2014gan}, defines a game-theoretic architecture where two neural networks—the Generator ($G$) and the Discriminator ($D$)—engage in a zero-sum competition. The generator $G$ is designed to map a latent noise vector $\boldsymbol{z}$ from a prior distribution $p_{\boldsymbol{z}}(\boldsymbol{z})$ to the data space, aiming to yield synthetic samples $G(\boldsymbol{z})$ that accurately mimic the real data distribution $p_{\text{data}}(\boldsymbol{x})$. Simultaneously, the discriminator $D$ acts as a binary classifier that is trained to distinguish between authentic samples $\boldsymbol{x}$ from the training set and synthesized samples $G(\boldsymbol{z})$ produced by the generator.

The fundamental objective of this adversarial process is modeled as a minimax game. Specifically, $D$ is optimized to maximize the probability of assigning the correct label to both real and fake data, while $G$ is concurrently trained to minimize the probability that $D$ correctly identifies its generated outputs as fake. The minimax objective function of a standard GAN is defined by the binary cross-entropy loss:
\begin{equation}
\min_{G} \max_{D} V(D, G) = \mathbb{E}_{\boldsymbol{x} \sim p_{\text{data}}(\boldsymbol{x})}[\log D(\boldsymbol{x})] + \mathbb{E}_{\boldsymbol{z} \sim p_{\boldsymbol{z}}(\boldsymbol{z})}[\log(1 - D(G(\boldsymbol{z})))]
\end{equation}
During the iterative training procedure, $D$ attempts to maximize $V(D, G)$ to accurately classify domain boundaries, while $G$ attempts to minimize $\log(1 - D(G(\boldsymbol{z})))$ to "fool" the discriminator into classifying synthetic inputs as real. In the specific context of industrial bearing fault diagnosis, GANs possess a powerful capability to learn the intricate, non-linear distributions of high-dimensional vibration signals \cite{radford2015dcgan}. However, standard architectures often suffer from mode collapse, where the generator produces a limited variety of samples, and significant training instability when applied to non-stationary mechanical time-series data \cite{li2021wgan_gp_fault, fan2021enhanced_gan}. Integrating such frameworks requires careful regularization and architectural considerations to ensure the generated signals maintain physical fidelity to actual mechanical behaviors \cite{liu2024gan_imbalanced_review}.

\subsection{Cycle-Consistent GAN Mechanism}
To overcome the significant limitation of requiring paired training datasets—which are often impossible to obtain in real-world mechanical fault scenarios—Zhu et al. \cite{zhu2017cyclegan} proposed the Cycle-Consistent GAN (CycleGAN). This architecture facilitates image-to-image or signal-to-signal translation between two unstructured domains $X$ and $Y$ without necessitating direct one-to-one correspondences. In our study, domain $X$ represents healthy vibration signals from a normal operational state, while domain $Y$ represents a collection of specific fault categories. The framework consists of two mapping functions (generators): $G: X \to Y$, which transforms normal signals into faulty ones, and $F: Y \to X$, which maps faulty signals back to the healthy domain. These are supported by two associated discriminators, $D_Y$ and $D_X$, which verify the authenticity of signals in their respective target domains.

The cornerstone of this architecture is the cycle-consistency loss, a structural constraint ensuring that translation from one domain to another followed by the inverse translation should reconstruct the original input exactly. This prevents the networks from simply mapping all inputs in domain $X$ to a single output in domain $Y$. Mathematically, the cycle-consistency loss is formulated as:
\begin{equation}
\mathcal{L}_{\text{cyc}}(G, F) = \mathbb{E}_{\boldsymbol{x} \sim p_{\text{data}}(\boldsymbol{x})} [\|F(G(\boldsymbol{x})) - \boldsymbol{x}\|_1] + \mathbb{E}_{\boldsymbol{y} \sim p_{\text{data}}(\boldsymbol{y})} [\|G(F(\boldsymbol{y})) - \boldsymbol{y}\|_1]
\end{equation}
By combining this cycle loss with the adversarial losses for both domains, the model learns to capture the underlying structural characteristics of vibration waveforms while performing cross-domain mapping. This mechanism is particularly advantageous for bearing fault diagnosis as it allows the generator to leverage the abundant temporal and spectral features present in the large-scale normal datasets to synthesize high-fidelity fault samples. Consequently, it effectively addresses the extreme data scarcity in minority fault classes by transferring structural knowledge from the health domain to the failure domain \cite{fernandez2022cyclegan_bearing, alotaibi2021cyclegan_signal, zhu2017cyclegan}.

\subsection{WGAN-GP and Lipschitz Constraints}
Standard GANs typically utilize the Jensen-Shannon (JS) divergence as their underlying distance metric. In many practical scenarios, if the distributions of real and generated data do not have a significant overlap, the JS divergence becomes constant, leading to the vanishing gradient problem and halting the learning process. To achieve stable training and provide a continuous gradient even when distributions are disjoint, Wasserstein GAN (WGAN) \cite{arjovsky2017wgan} utilizes the Earth Mover's (EM) distance (also known as the Wasserstein-1 distance) as the optimization target. The Wasserstein distance between the real distribution $\mathbb{P}_r$ and the generated distribution $\mathbb{P}_g$ is defined via the Kantorovich-Rubinstein duality as:
\begin{equation}
W(\mathbb{P}_r, \mathbb{P}_g) = \frac{1}{K} \sup_{\|f\|_L \le K} \mathbb{E}_{\boldsymbol{x} \sim \mathbb{P}_r}[f(\boldsymbol{x})] - \mathbb{E}_{\tilde{\boldsymbol{x}} \sim \mathbb{P}_g}[f(\tilde{\boldsymbol{x}})]
\end{equation}
where $f$ is a 1-Lipschitz continuous function. To strictly enforce this Lipschitz constraint without the detrimental side effects of weight clipping (such as gradient capacity reduction or parameter saturation), the WGAN with Gradient Penalty (WGAN-GP) \cite{gulrajani2017wgangp} was introduced. The GP term serves as a soft constraint that penalizes the gradient norm of the discriminator (critic) output with respect to its input, forcing it toward unity:
\begin{equation}
\mathcal{L}_{\text{GP}} = \lambda \mathbb{E}_{\hat{\boldsymbol{x}} \sim p_{\hat{\boldsymbol{x}}}} [(\|\nabla_{\hat{\boldsymbol{x}}} D(\hat{\boldsymbol{x}})\|_2 - 1)^2]
\end{equation}
In this expression, $\hat{\boldsymbol{x}}$ is sampled uniformly along straight lines between pairs of points from the real data distribution and the generator distribution, and $\lambda$ is a penalty coefficient. Integrating the WGP strategy into our CycleGAN-based framework effectively stabilizes the dual-mapping training process, preventing gradient explosions during the long-sequence signal generation. This ensures that the synthesized faults maintain physical consistency with the original mechanical signatures through rigorous mathematical regularization \cite{han2020wgangp_bearing, meng2022iacgan, gulrajani2017wgangp}.
